% ═════════════════════════════════════════════════════════════════════════════
%  PREAMBLE
% ═════════════════════════════════════════════════════════════════════════════

% ── Encoding and fonts ───────────────────────────────────────────────────────
\usepackage[T1]{fontenc}
\usepackage[utf8]{inputenc}
\usepackage{lmodern}
\usepackage[english]{babel}

% ── Page geometry ────────────────────────────────────────────────────────────
% Conventional thesis page: 11pt body, margins near 3 cm with a little extra
% on the binding edge, giving a measure of roughly 85 characters.
\usepackage[a4paper, inner=2.9cm, outer=2.7cm, top=2.5cm, bottom=2.7cm,
            headheight=14pt, headsep=13pt]{geometry}
\usepackage{setspace}
% 1.3 keeps the lines clearly separated at this measure while fitting
% appreciably more text per page than one-and-a-half spacing.
\setstretch{1.23}

% ── Mathematics ──────────────────────────────────────────────────────────────
\usepackage{amsmath, amssymb, amsthm}
% Number equations by section (2.1, 2.2, ...) rather than continuously, so a
% reference locates the equation as well as naming it.
\numberwithin{equation}{section}

% ── Tables and figures ───────────────────────────────────────────────────────
\usepackage{graphicx}
% Figures have one source of truth: analysis/results/figures/. The relative path keeps the
% paper reproducible without maintaining a second, potentially stale, copy.
\graphicspath{{../analysis/results/figures/}{assets/}{./}}
\usepackage{booktabs}
\usepackage{array}
% Booktabs rules need room to breathe; the default row height crowds them.
\renewcommand{\arraystretch}{1.18}
\usepackage{float}
\usepackage[font=small, labelfont=bf, labelsep=period,
            width=0.94\textwidth, justification=justified,
            singlelinecheck=true]{caption}

% ── Lists, headers, references ───────────────────────────────────────────────
\usepackage{enumitem}
\setlist[enumerate]{itemsep=2pt, topsep=4pt}
\setlist[itemize]{itemsep=2pt, topsep=4pt}
\usepackage{natbib}
\bibliographystyle{plainnat}
\usepackage{fancyhdr}
\usepackage{titlesec}
\usepackage{tocloft}
% Open up the contents list: entries are few, so they can afford the space.
\setlength{\cftbeforesecskip}{0.7em}
\setlength{\cftbeforesubsecskip}{0.25em}
\renewcommand{\cftsecfont}{\bfseries}
\renewcommand{\cftsecpagefont}{\bfseries}
\renewcommand{\cftsecleader}{\cftdotfill{\cftdotsep}}
\renewcommand{\cftsubsecleader}{\hfill}
\renewcommand{\cftsubsecafterpnum}{\hspace*{0pt}}
\usepackage{microtype}

% ── Hyperlinks (loaded late) ─────────────────────────────────────────────────
\usepackage{xcolor}
\definecolor{linknavy}{RGB}{20,60,110}
\definecolor{citeink}{RGB}{55,70,95}
\usepackage[hidelinks]{hyperref}
% Long dataset URLs cannot be hyphenated; let them break at any character so a
% single reference does not force loose lines through the whole entry.
\usepackage{url}
\def\UrlBreaks{\do\/\do-\do.\do:\do_\do?\do\&\do=\do+}
\Urlmuskip=0mu plus 2mu
\hypersetup{
  colorlinks = true,
  linkcolor  = citeink,
  citecolor  = citeink,
  urlcolor   = linknavy,
  pdftitle   = {The Shape of Evidence: Fuzzy Prediction Sets for Downstream Decisions},
  pdfauthor  = {Antonis Konstantinidis},
  bookmarksnumbered = true
}

% ═════════════════════════════════════════════════════════════════════════════
%  STYLING
% ═════════════════════════════════════════════════════════════════════════════

% Running header: section name left, page number right, set in small caps so
% it recedes behind the text rather than competing with the headings.
\pagestyle{fancy}
\fancyhf{}
\renewcommand{\headrulewidth}{0.4pt}
\fancyhead[L]{\footnotesize\scshape\nouppercase{\leftmark}}
\fancyhead[R]{\footnotesize\thepage}
\fancypagestyle{plain}{\fancyhf{}\renewcommand{\headrulewidth}{0pt}%
  \fancyfoot[C]{\footnotesize\thepage}}

% Section headings. Generous space above, tighter below, so each heading binds
% to the text it introduces rather than floating between blocks.
\titleformat{\section}{\normalfont\Large\bfseries}{\thesection}{0.7em}{}
\titleformat{\subsection}{\normalfont\large\bfseries}{\thesubsection}{0.7em}{}
% Every major section starts on a new page. Subsections remain attached to the
% surrounding text because they mark divisions within a section, not new
% chapters.
\titlespacing*{\section}{0pt}{3.4ex plus 1ex minus .3ex}{1.5ex plus .2ex}
\titlespacing*{\subsection}{0pt}{2.4ex plus .8ex minus .2ex}{0.9ex plus .2ex}
\newcommand{\sectionbreak}{\clearpage}

% Floats: keep them clear of the surrounding text.
\setlength{\textfloatsep}{14pt plus 3pt minus 3pt}
\setlength{\intextsep}{12pt plus 3pt minus 3pt}
\setlength{\abovecaptionskip}{7pt}
\setlength{\belowcaptionskip}{2pt}

% Display equations and the operators inside them get more air. Size commands
% reset the display skips, so they are reapplied whenever \normalsize runs.
\newcommand{\displayair}{%
  \setlength{\abovedisplayskip}{13pt plus 3pt minus 2pt}%
  \setlength{\belowdisplayskip}{13pt plus 3pt minus 2pt}%
  \setlength{\abovedisplayshortskip}{8pt plus 3pt}%
  \setlength{\belowdisplayshortskip}{10pt plus 3pt}}
\usepackage{etoolbox}
\appto{\normalsize}{\displayair}
\AtBeginDocument{\displayair}
\setlength{\jot}{6pt}
\medmuskip=4.5mu plus 2mu minus 2mu
\thickmuskip=6mu plus 4mu minus 1mu

% Discourage widows and orphans.
\widowpenalty=10000
\clubpenalty=10000

% ═════════════════════════════════════════════════════════════════════════════
%  THEOREM ENVIRONMENTS
% ═════════════════════════════════════════════════════════════════════════════
% Theorem-like results are separated from the prose that frames them.
\newlength{\thmsep}\setlength{\thmsep}{1.1\baselineskip}
\makeatletter
\def\thm@space@setup{%
  \thm@preskip=\thmsep \thm@postskip=\thmsep}
\makeatother
\theoremstyle{plain}
\newtheorem{theorem}{Theorem}[section]
\newtheorem{proposition}[theorem]{Proposition}
\newtheorem{lemma}[theorem]{Lemma}
\newtheorem{corollary}[theorem]{Corollary}
\theoremstyle{definition}
\newtheorem{definition}[theorem]{Definition}
\theoremstyle{remark}
\newtheorem{remark}[theorem]{Remark}

% ═════════════════════════════════════════════════════════════════════════════
%  MATH SHORTHAND
% ═════════════════════════════════════════════════════════════════════════════
\DeclareMathOperator*{\argmin}{arg\,min}

